\documentclass[11pt,a4paper]{article}
\usepackage[margin=25mm]{geometry}
\usepackage[T1]{fontenc}
\usepackage{lmodern}
\usepackage{booktabs,tabularx,array}
\usepackage{xcolor}
\usepackage{hyperref}
\hypersetup{colorlinks=true,linkcolor=blue,urlcolor=blue}
\title{Cryptography and Network Security\\Integrated Situation Report}
\begin{document}
\maketitle

\begin{abstract}
This report documents a risk assessment, an AES-256-GCM file encryption and decryption utility, SHA-256 integrity checking, and a routed nftables policy template for an authorised laboratory. Network identifiers and service details were not supplied in the assignment, so they remain marked placeholders. Laboratory firewall tests have not been executed; expected outcomes are distinguished from actual evidence.
\end{abstract}

\section{Scope and assets}
The scenario concerns student records on a central server and file transfers between two campuses. It identifies weak staff passwords, outdated software, unencrypted transfers, guest reachability to the records server, and repeated attempts from an unfamiliar external address. No real student data or production network was used in preparing this project.

\begin{tabularx}{\textwidth}{>{\raggedright\arraybackslash}p{0.23\textwidth}XX}
\toprule
Asset & Vulnerability or threat & Possible consequence\\
\midrule
Student records and server & Guest access and unfamiliar external attempts & Unauthorised disclosure, alteration, loss, or service disruption\\
Staff accounts and systems & Weak passwords and outdated software & Account compromise, exploitation, malware, unauthorised access\\
Inter-campus files and transfer service & Unencrypted transfers & Interception or modification in transit\\
\bottomrule
\end{tabularx}

\section{Risk ranking and controls}
Ratings are qualitative estimates based only on the scenario, not measured incident data.
\begin{tabularx}{\textwidth}{c>{\raggedright\arraybackslash}p{0.27\textwidth}ccX}
\toprule
Rank & Risk & Likelihood & Impact & Reason and recommended control\\
\midrule
1 & Guest or external access to records server & High & High & Guest reachability and repeated external attempts are stated. Apply default-deny filtering, block guest and other inbound sources, allow only authorised staff to the assessor-selected service, and review logs.\\
2 & Staff account compromise or software exploitation & Medium & High & Weak passwords and outdated software are confirmed, but no compromise is reported. Enforce unique strong passwords and MFA where available, patch supported systems, minimise privileges, and review access.\\
3 & Interception or tampering in transit & Medium & High & Transfers are unencrypted, though no interception is reported. Use institution-approved authenticated encryption for transfer and verify integrity over a trusted channel.\\
\bottomrule
\end{tabularx}

\section{Encryption and integrity toolkit}
The Python CLI uses AES-256-GCM from the \texttt{cryptography} package. A fresh 96-bit nonce is generated for each encryption. The output has a small format header followed by the nonce and authenticated ciphertext; decryption rejects malformed data, wrong keys, and altered ciphertext. Inputs and outputs must be different files, and the tool refuses to overwrite outputs. A random 256-bit key is generated separately and saved outside the repository with restricted permissions where the platform supports them.

The \texttt{hash} command calculates SHA-256 and can save a JSON baseline. The \texttt{verify} command compares a later file digest with that baseline and reports match or mismatch. The baseline must be protected independently: an unkeyed hash alone does not establish who created it.

\subsection{Reproducible evidence}
The synthetic checks run for this project are recorded in \texttt{crypto\_tests.md}: key generation, encryption, decryption with byte-for-byte round trip, matching and changed-file SHA-256 verification, missing input, and malformed ciphertext. The changed-file check returned exit status 1 as expected; malformed inputs returned concise errors with exit status 2. The CLI also handles malformed key files, invalid hash records, and overwrite attempts with a concise error rather than an unhandled exception.

\section{Firewall policy and connection testing}
The included \texttt{firewall.nft.template} is a routed Linux nftables example with a default-drop forwarding policy. It blocks guest traffic to the records server before accepting established traffic, allows the staff CIDR to the server only on the assessor-approved service port, denies other sources to that service, and drops remaining forwarded packets.

The assignment does not provide the server address, staff/guest CIDRs, gateway interface, or approved service, transport protocol, and port. The template therefore cannot safely be applied until the authorised administrator supplies and reviews those values. Confirm that the approved protocol is supported by the nftables transport-header expression in the lab version; adapt the rules if it is not, if the records server is not routed through this gateway, or if the lab uses another firewall platform.

\begin{tabularx}{\textwidth}{>{\raggedright\arraybackslash}p{0.18\textwidth}XX}
\toprule
Test & Expected outcome & Actual outcome\\
\midrule
Staff client to approved service & Permitted while service is listening & Not executed; requires authorised lab and supplied values\\
Guest client to approved service & Blocked & Not executed; requires authorised lab and supplied values\\
Other source to approved service & Blocked & Not executed; requires authorised lab and supplied values\\
\bottomrule
\end{tabularx}

Detailed command templates and the safe execution procedure are in \texttt{filter\_tests.md}. No firewall rules were applied and no network tests were run for this project build. The expected results above are policy expectations, not observed results.

\section{Conclusion}
The project provides a structured assessment of the three scenario risks, an authenticated file encryption tool with SHA-256 baseline checks, and a default-deny firewall template with a reproducible three-case test plan. Completion of the live firewall portion requires the assessor's exact network and service values plus authorised lab access. The actual test outcomes and student-identifying information should be added before final submission.

GitHub repository URL: \texttt{[https://github.com/olivieryeada8-del/cryptography-network-security-exam.git]}.

\end{document}
